% ---------------------------------------------------------
% TABELA EMENTA 1: Artes — Ano 1
% ---------------------------------------------------------
\begin{xltabular}{\linewidth}{|X|p{2.5cm}|p{2.5cm}|}
\hline
\multirow{3}{=}{\textcolor{ifscgreen}{\textbf{Unidade Curricular:}}\\[3pt]\textbf{\large Artes — Ano 1}} & \multicolumn{2}{p{\dimexpr5cm+2\tabcolsep+\arrayrulewidth\relax}|}{\textcolor{ifscgreen}{\textbf{Semestre:}} \textbf{1º e 2º}} \\ \cline{2-3}
 & \textcolor{ifscgreen}{\textbf{CH EaD*:}} & \textcolor{ifscgreen}{\textbf{CH Total*:}} \\ \cline{2-3}
 & \textbf{00 h} & \textbf{80 h} \\ \hline
\endhead
\multicolumn{3}{|p{\dimexpr\linewidth-2\tabcolsep-2\arrayrulewidth\relax}|}{\textcolor{ifscgreen}{\textbf{Objetivos:}}} \\ \hline
\multicolumn{3}{|p{\dimexpr\linewidth-2\tabcolsep-2\arrayrulewidth\relax}|}{
\begin{itemize}[leftmargin=*,noitemsep,topsep=2pt]
 \item Compreender a Arte como construção histórica, social e cultural e suas imbricações.
 \item Refletir sobre as relações que envolvem os processos de construção e fruição da Arte.
 \item Expressar e comunicar ideias e sentimentos por meio de atividades de práticas corporais.
 \item Desenvolver o senso crítico em relação à arte, à sociedade, à tecnologia e ao meio ambiente.
 \item Interpretar e analisar o significado de obras de arte a partir de diferentes perspectivas.
 \item Identificar, experienciar e criar nas diferentes linguagens da arte do corpo em cena (dança, teatro e performance).
\end{itemize}
} \\ \hline
\multicolumn{3}{|p{\dimexpr\linewidth-2\tabcolsep-2\arrayrulewidth\relax}|}{\textcolor{ifscgreen}{\textbf{Conteúdos:}}} \\ \hline
\multicolumn{3}{|p{\dimexpr\linewidth-2\tabcolsep-2\arrayrulewidth\relax}|}{
\begin{itemize}[leftmargin=*,noitemsep,topsep=2pt]
 \item Conceitos básicos: linguagem, objeto de conhecimento, significação e produto.
 \item Arte contextualizada social, política, histórica e economicamente.
 \item Manifestações artísticas no decorrer do tempo e do espaço.
 \item Apreciação e análise de produções artísticas internacionais, nacionais e locais.
 \item Arte, ciência e meio ambiente: a não separação entre natureza e cultura.
 \item Arte do corpo em cena na atualidade: apreciação, análise e prática.
 \item O corpo como instrumento de expressão artística.
\end{itemize}
} \\ \hline
\multicolumn{3}{|p{\dimexpr\linewidth-2\tabcolsep-2\arrayrulewidth\relax}|}{\textcolor{ifscgreen}{\textbf{Estratégias de Ensino e Aprendizagem:}}} \\ \hline
\multicolumn{3}{|p{\dimexpr\linewidth-2\tabcolsep-2\arrayrulewidth\relax}|}{- **Metodologia:** Aulas participativas e dialogadas, partindo do conhecimento prévio do estudante, utilizando imagens, textos, vídeos, músicas e corpo; atividades focadas na indissociabilidade entre teoria e prática, com atividades corporais; desenvolvimento de trabalhos, exercícios, pesquisas, seminários e obras artísticas; atividades individuais e em grupos; jogos teatrais e exercícios de improvisação; práticas de dança contemporânea e danças tradicionais brasileiras; utilização da sala de artes para atividades práticas.
- **Avaliação:** Não especificado no documento fonte.} \\ \hline
\multicolumn{3}{|p{\dimexpr\linewidth-2\tabcolsep-2\arrayrulewidth\relax}|}{\textcolor{ifscgreen}{\textbf{Bibliografia Básica:}}} \\ \hline
\multicolumn{3}{|p{\dimexpr\linewidth-2\tabcolsep-2\arrayrulewidth\relax}|}{1. GOMBRICH, E. H. \textbf{A história da arte}. Tradução de Álvaro Cabral. 16. ed. Rio de Janeiro: Livros Técnicos e Científicos, 1999.
2. PROENÇA, G. \textbf{História da arte}. 17. ed. São Paulo: Ática, 2010.} \\ \hline
\multicolumn{3}{|p{\dimexpr\linewidth-2\tabcolsep-2\arrayrulewidth\relax}|}{\textcolor{ifscgreen}{\textbf{Bibliografia Complementar:}}} \\ \hline
\multicolumn{3}{|p{\dimexpr\linewidth-2\tabcolsep-2\arrayrulewidth\relax}|}{1. AGRA, Lucio. \textbf{História da arte do século XX: idéias e movimentos}. 2. ed. rev. atual. São Paulo: Anhembi, 2004.
2. CERTEAU, Michel de. \textbf{A invenção do cotidiano: artes de fazer}. Tradução de Ephraim Ferreira Alves. 16. ed. Petrópolis, RJ: Vozes, 2009.
3. LARAIA, R. de B. \textbf{Cultura: um conceito antropológico}. 26. reimp. Rio de Janeiro: Jorge Zahar, 1986.
4. SOUZA, Ana Lúcia Silva. \textbf{Letramentos de reexistência: poesia, grafite, música, dança: hip-hop}. São Paulo: Parábola Editorial, 2011.} \\ \hline
\end{xltabular}

\vspace{0.4cm}


% ---------------------------------------------------------
% TABELA EMENTA 2: Artes — Ano 2
% ---------------------------------------------------------
\begin{xltabular}{\linewidth}{|X|p{2.5cm}|p{2.5cm}|}
\hline
\multirow{3}{=}{\textcolor{ifscgreen}{\textbf{Unidade Curricular:}}\\[3pt]\textbf{\large Artes — Ano 2}} & \multicolumn{2}{p{\dimexpr5cm+2\tabcolsep+\arrayrulewidth\relax}|}{\textcolor{ifscgreen}{\textbf{Semestre:}} \textbf{3º e 4º}} \\ \cline{2-3}
 & \textcolor{ifscgreen}{\textbf{CH EaD*:}} & \textcolor{ifscgreen}{\textbf{CH Total*:}} \\ \cline{2-3}
 & \textbf{00 h} & \textbf{80 h} \\ \hline
\endhead
\multicolumn{3}{|p{\dimexpr\linewidth-2\tabcolsep-2\arrayrulewidth\relax}|}{\textcolor{ifscgreen}{\textbf{Objetivos:}}} \\ \hline
\multicolumn{3}{|p{\dimexpr\linewidth-2\tabcolsep-2\arrayrulewidth\relax}|}{
\begin{itemize}[leftmargin=*,noitemsep,topsep=2pt]
 \item Refletir sobre os conceitos básicos dos fenômenos artísticos ocidentais e suas relações com outras formas de pensamento.
 \item Analisar o uso de manifestações tradicionais brasileiras como constituintes de saberes, territórios e identidades.
 \item Vivenciar práticas corporais tradicionais.
 \item Experienciar e criar nas diferentes linguagens da arte do corpo em cena (dança, teatro e performance).
 \item Valorizar a diversidade cultural brasileira.
 \item Analisar o papel da arte como ferramenta de denúncia e conscientização.
 \item Refletir sobre a relação entre arte, política, ciência, tecnologia e meio ambiente.
 \item Estabelecer diálogos entre a arte contemporânea e os saberes tradicionais.
\end{itemize}
} \\ \hline
\multicolumn{3}{|p{\dimexpr\linewidth-2\tabcolsep-2\arrayrulewidth\relax}|}{\textcolor{ifscgreen}{\textbf{Conteúdos:}}} \\ \hline
\multicolumn{3}{|p{\dimexpr\linewidth-2\tabcolsep-2\arrayrulewidth\relax}|}{
\begin{itemize}[leftmargin=*,noitemsep,topsep=2pt]
 \item Manifestações populares e tradicionais no Brasil.
 \item Cultura indígena, afro-brasileira e africana.
 \item Artivismo.
 \item Abordagens em arte contra a hegemonia, em diálogo com os saberes e as tradições decoloniais.
 \item Práticas de movimentos, passos e técnicas em dança e teatro.
 \item Estudo teórico-prático de técnicas de improvisação em dança e teatro.
 \item Panorama histórico da dança e do teatro no Brasil.
 \item Pesquisa, elaboração de ideias, criação, produção e apresentação de obras em artes do corpo em cena.
 \item Improvisação e criação coletiva.
\end{itemize}
} \\ \hline
\multicolumn{3}{|p{\dimexpr\linewidth-2\tabcolsep-2\arrayrulewidth\relax}|}{\textcolor{ifscgreen}{\textbf{Estratégias de Ensino e Aprendizagem:}}} \\ \hline
\multicolumn{3}{|p{\dimexpr\linewidth-2\tabcolsep-2\arrayrulewidth\relax}|}{- **Metodologia:** Aulas participativas e dialogadas, partindo do conhecimento prévio do estudante, utilizando imagens, textos, vídeos, músicas e corpo; atividades focadas na indissociabilidade entre teoria e prática, com atividades corporais; desenvolvimento de trabalhos, exercícios, pesquisas, seminários e obras artísticas; atividades individuais e em grupos; jogos teatrais e exercícios de improvisação; práticas de dança contemporânea e danças tradicionais brasileiras; utilização da sala de artes para atividades práticas.
- **Avaliação:** Não especificado no documento fonte.} \\ \hline
\multicolumn{3}{|p{\dimexpr\linewidth-2\tabcolsep-2\arrayrulewidth\relax}|}{\textcolor{ifscgreen}{\textbf{Bibliografia Básica:}}} \\ \hline
\multicolumn{3}{|p{\dimexpr\linewidth-2\tabcolsep-2\arrayrulewidth\relax}|}{1. CONDURU, Roberto. \textbf{Arte afro-brasileira}. Belo Horizonte: C/Arte, 2007.
2. GOMBRICH, E. H. \textbf{A história da arte}. Tradução de Álvaro Cabral. 16. ed. Rio de Janeiro: Livros Técnicos e Científicos, 1999.
3. PROENÇA, G. \textbf{História da arte}. 17. ed. São Paulo: Ática, 2010.} \\ \hline
\multicolumn{3}{|p{\dimexpr\linewidth-2\tabcolsep-2\arrayrulewidth\relax}|}{\textcolor{ifscgreen}{\textbf{Bibliografia Complementar:}}} \\ \hline
\multicolumn{3}{|p{\dimexpr\linewidth-2\tabcolsep-2\arrayrulewidth\relax}|}{1. NILSON, Afonso. \textbf{Seis textos breves para estudantes de teatro}. Florianópolis: Letras Contemporâneas, 2017.
2. MACHADO, Lúcia. \textbf{A modernidade no teatro: [ali e aqui]: reflexos estilhaçados}. Recife: Ed. do autor, 2009.
3. PIMENTEL, Spency. \textbf{O índio que mora na nossa cabeça: sobre as dificuldades para entender os povos indígenas}. São Paulo: Prumo, 2012.
4. SANT'ANNA, Márcia (org.). \textbf{Os sambas brasileiros: diversidade, apropriação e salvaguarda}. Brasília, DF: IPHAN, 2011.} \\ \hline
\end{xltabular}

\vspace{0.4cm}

